%% kapa: exam=2 round=33 year=2022 period=1 subject=appraisal_practice subject_ko=감정평가실무 seq=4 number=4 points=10
%% source: Q-Net 감정평가사 기출문제 (cst00309) artlSeq=5210838 file=raw/1008/5210838/2022년 제33회 감정평가사 2차시험 1교시 문제지(감정평가실무).pdf
%% conversion: kordoc 4.24.51 -> tools/gen2.py (auto); review=proofread-gpt-astra; flags=none
\begin{longquestion}{【문제 4】}
부동산을 명도받기 위한 소송을 제기한 임대인 원고에 맞서, 임차인이자 개인사업자인 피고는 「상가건물 임대차보호법」의 권리금 회수기회 보호 등의 규정을 들어 원고로부터 권리금의 지급을 요청하는 ‘감정신청서’를 법원에 제출하였다. 재판장은 피고의 감정신청을 받아들여 감정평가사 甲에게 권리금에 대한 감정평가를 의뢰하였다. 제시된 자료를 활용하여 대상 사업체의 권리금을 산정하시오. (단, 권리금은 시설권리금, 영업권리금, 바닥권리금으로 구분하여 제시할 것) (10점) \par
< 자료 1 > 기본적 사항 \par
1. 기준가치: 시장가치 \par
2. 기준시점: 2022년 7월 16일 \par
3. 대상 사업체의 개황 \par
1) 소재지: C시 D동 120 \par
2) 업종: 커피숍 \par
3) 개업일: 2017년 1월 1일 \par
4) 면적: 120m\textsuperscript{2} \par
< 자료 2 > 시설권리금 자료 \par
1. 시설권리금 대상인 유형재산은 인테리어뿐이며, 사업자는 개업일 당시 인테리어 비용 600,000원/m\textsuperscript{2} 소요되었다는 자료를 제출하였고, 제반 상황을 고려할 때 비용은 적정한 것으로 판단됨 \par
2. 기준시점의 재조달원가는 개업일 당시 비용에 건축공사비지수를 적용하여 산정하며, 조사된 건축공사비지수는 다음 표와 같다. \par
{\small\centering\renewcommand{\arraystretch}{1.25}
\resizebox{\ifdim\width>\linewidth\linewidth\else\width\fi}{!}{
\begin{tabular}{|c|c|c|}
\hline
\makecell{구분} & \makecell{2017년 1월} & \makecell{2022년 7월} \\
\hline
\makecell{건축공사비지수} & \makecell{112} & \makecell{147} \\
\hline
\end{tabular}}\par}

- 건축공사비 변동률 산정은 일할계산하지 않고, 해당 월에 고시된 건축공사비지수를 적용하며, 소수점 넷째 자리를 반올림하여 셋째 자리까지 표기함 \par
3. 단가 산정은 천원 단위 미만에서 반올림하여 천원 단위까지 표기함 \par
4. 감가수정은 정액법에 따르고, 총내용연수는 동종 업의 인테리어 수명 주기를 고려하여 10년으로 하며, 연단위 만년감가를 적용함 \par
< 자료 3 > 영업권리금 자료 \par
1. 영업권리금 산정을 위한 영업이익은 기준시점 이전 3년의 평균영업이익인 연간 23,000,000원으로 하였음 \par
2. 개인사업자로서 영업이익에서 공제해야하는 비용은 자가인건비 상당액으로 기준시점 이전 3년 평균인 연간 19,000,000원으로 하였으며, 감가상각비는 고려하지 않음 \par
3. 무형재산 귀속 영업이익은 브랜드를 선호하는 업종의 특성을 고려할 때, 50\%를 적용하는 것이 타당한 것으로 판단됨 \par
4. 인근지역 브랜드 커피숍 증가로 기준시점 이후 영업이익은 동일할 것으로 추정함 \par
5. 할인기간은 5년으로 하고, 기준시점 이후 5년간 추정된 영업이익에 대응하는 할인율은 아래 표와 같음 \par
{\small\centering\renewcommand{\arraystretch}{1.25}
\resizebox{\ifdim\width>\linewidth\linewidth\else\width\fi}{!}{
\begin{tabular}{|c|c|c|c|c|c|}
\hline
\makecell{구분} & \makecell{1년} & \makecell{2년} & \makecell{3년} & \makecell{4년} & \makecell{5년} \\
\hline
\makecell{할인율} & \makecell{0.899} & \makecell{0.808} & \makecell{0.726} & \makecell{0.653} & \makecell{0.587} \\
\hline
\end{tabular}}\par}

< 자료 4 > 바닥권리금 자료 \par
대상 사업체가 속한 상권은 위치와 업종, 가로의 상태에 따라 일부 바닥권리금이 형성되는 상가가 있으나, 시설권리금과 영업권리금을 받을 수 있는 상가는 별도의 바닥권리금이 없는 것으로 조사됨 \par
\end{longquestion}
